\section{Affiliated rationality and regular fractions}\label{sec:fractions}
All tensor products in this section that carry no bar are algebraic over the indicated field.
Set $U:=L^0(\M)$.
Under $\Delta(X)=d$, the algebra $\mathcal D_X$ is a division ring, the image of the free field under evaluation \cite[Theorem~1.1]{MSY23}.
Define
\[
 \Psi_i(g\otimes h):=gX_i\otimes h-g\otimes X_ih,
 \qquad g,h\in U.
\]
The maps $\Psi_i$ commute with the outer $U$-bimodule actions.

\begin{proposition}\label{prop:intrinsic}
For an arbitrary bounded tuple with $\Delta(X)=d$ and every $u\in U$,
\begin{equation}\label{eq:affiliated-tensor}
 u\in\mathcal D_X
 \quad\Longleftrightarrow\quad
 u\otimes1-1\otimes u=\sum_i\Psi_i(k_i),\qquad k_i\in U\otimes_\C U.
\end{equation}
The tensors $k_i$ are unique when they exist.
Moreover \eqref{eq:affiliated-tensor} is equivalent to existence of one bounded regular fraction and finite-rank operators satisfying \eqref{eq:boundary}.
At a fixed fraction the solution tuple is unique.
\end{proposition}
\begin{proof}
For a scalar rational expression defining $u\in\mathcal D_X$, form its derivative tensors by the product rule and
\[
 \partial_i(v^{-1})=-v^{-1}\cdot(\partial_i v)\cdot v^{-1}.
\]
The map $u\mapsto u\otimes1-1\otimes u$ has the same product and inverse rules and takes $X_j$ to $\Psi_j(1\otimes1)$.
Induction proves the forward direction.
For the reverse, project to $U\otimes_{\mathcal D_X}U$.
Balancing kills every $\Psi_i(k_i)$, so $u\otimes_{\mathcal D_X}1=1\otimes_{\mathcal D_X}u$.
If $u\notin\mathcal D_X$, the vectors $1,u$ are right-linearly independent over this division ring.
Extend them to a right basis and choose a right-linear functional taking $1$ to zero and $u$ to one.
Applying this functional on the first leg contradicts the equality.

For uniqueness, clear the finitely many first tensor legs on the left and second legs on the right by bounded regular $s,t$ so that $sk_it\in\M\otimes_\C\M$.
One may obtain these denominators from an affiliated row $G$ and column $B$ as $(1+GG^*)^{-1/2}$ and $(1+B^*B)^{-1/2}$.
For algebraic $L^2$ tensors,
\begin{equation}\label{eq:psi-theta}
 [\rho(X_i),\Theta(k)]=\Theta(\Psi_i(k)).
\end{equation}
Thus a zero sum of the $\Psi_i(k_i)$ gives a finite-rank cycle after clearing.
Equation~\eqref{eq:cycles} makes every cleared tensor zero, and invertibility of $s,t$ in $U$ gives uniqueness.

For the fraction assertion, clear $u$ as well as all the tensor legs in \eqref{eq:affiliated-tensor}, put $a:=su$, $b:=ut$, and take $T_i:=\Theta(sk_it)$.
Then
\[
 \sum_i\Psi_i(sk_it)=a\otimes t-s\otimes b,
\]
which gives \eqref{eq:boundary} by \eqref{eq:psi-theta}.
Conversely, represent each finite-rank $T_i$ by an algebraic $L^2$ tensor, include those tensors in $U\otimes_\C U$, and divide the resulting identity on the left by $s$ and on the right by $t$.
This gives \eqref{eq:affiliated-tensor}.
Uniqueness at fixed $Q$ also follows by subtracting two solutions in \eqref{eq:boundary} and using \eqref{eq:cycles}.
\end{proof}

The next argument explains why the specified duals realize the certificate at every fraction under the stronger domain hypotheses.
We give the algebraic fraction calculations before using normality.
Let $E_i$ be operators defined on $\M\Omega$, with $\M\Omega\subseteq\Dom E_i^*$ and $[E_i,X_j]=\delta_{ij}P_\Omega$ there.
For a fraction $Q$ put
\[
 K_i^Qz:=sE_i(bz)-aE_i(tz),\qquad z\in\M\Omega.
\]
Its adjoint is defined on $\M\Omega$ by testing against
\begin{equation}\label{eq:fraction-adjoint}
 b^*E_i^*(s^*w)-t^*E_i^*(a^*w).
\end{equation}
Consequently $K_i^Q$ is closable.
A densely defined closable finite-rank operator has a bounded finite-rank closure of the same rank \cite[Problem~III.5.18]{Kat95}.

For two fractions of the same $u$, choose bounded regular common left and right multiples
\[
 xs=ys_1,\qquad tx'=t_1y'.
\]
Such multiples exist by expressing $ss_1^{-1}$ as a bounded left fraction and $t^{-1}t_1$ as a bounded right fraction; the auxiliary numerators are regular because these ratios are units of $U$.
The fraction identities give $xa=ya_1$, $bx'=b_1y'$, whence
\begin{equation}\label{eq:common-defect}
 yK_i^{Q_1}(y'z)=xK_i^Q(x'z)\qquad(z\in\M\Omega).
\end{equation}
Finite rank can be cancelled through a bounded injective left multiplier and a bounded regular right input multiplier.
Indeed, the latter sends $\M\Omega$ to a dense subspace of $H$; the restriction of a closable operator to this subspace, if finite rank, has a bounded closure on all $H$.
Its graph is contained in the original closed graph, forcing the same bounded extension and finite rank for the original restriction.
The numerical ranks agree by density and injectivity.
Thus finite rank and its value in each coordinate are independent of the chosen fraction.

The class of elements whose $K_i^Q$ are all finite rank is a unital complex algebra closed under affiliated inverses.
Addition uses common left and right denominators and adds the defects.
For multiplication, write $u_l=s_l^{-1}a_l=b_lt_l^{-1}$ and choose
\[
 a_1s_2^{-1}=x^{-1}y,\qquad t_1^{-1}b_2=pq^{-1},
\]
with $x,q$ regular; no regularity of $y,p$ is needed.
A fraction for $u_1u_2$ is $(xs_1,ya_2,b_1p,t_2q)$ and its defect is
\[
 xK_i^{Q_1}(pz)+yK_i^{Q_2}(qz).
\]
The two middle terms cancel because $xa_1=ys_2$ and $t_1p=b_2q$.
For an invertible $u$, the inverse fraction $(a,s,t,b)$ has defect $-K_i^Q$.
Finally the generator fractions have defects $[E_i,X_j]=\delta_{ij}P_\Omega$.
We have proved that every $u\in\mathcal D_X$ has finite-rank defects at every bounded regular fraction.

\begin{lemma}\label{lem:fundamental}
Suppose the $X_j$ are bounded normal operators generating $\M$ and the $E_i$ satisfy the domain and commutator hypotheses just used.
If the $K_i^Q$ have bounded extensions, then these extensions satisfy \eqref{eq:boundary}.
\end{lemma}
\begin{proof}
The closed graph theorem for their closures restricted to $\M$ makes the $E_i$ bounded $\M\to H$.
Their densely defined adjoints imply weak continuity on operator-norm-bounded sets that converge in $L^2$: test first against the adjoint domain and then use boundedness and density.
Define on $\M\Omega$
\[
 C_0:=\sum_i(\rho(X_i)E_i-E_i\rho(X_i))+P_\Omega.
\]
This map is bounded $\M\to H$ and has the same bounded-set weak continuity.
The generator commutators and traciality give $C_0X_j=X_jC_0$.
The homogeneous argument in Lemma~\ref{lem:commutation} gives commutation with every $X_j^*$.
Kaplansky approximation and bounded-set weak continuity therefore give $C_0(az)=aC_0z$ for all $a,z\in\M$.
Since $sb=at$, direct expansion yields
\[
 \sum_i[\rho(X_i),K_i^Q]z
 =s(C_0-P_\Omega)(bz)-a(C_0-P_\Omega)(tz)
 =(aP_\Omega t-sP_\Omega b)z.
\]
Both sides are bounded and agree on the dense subspace $\M\Omega$, proving the identity.
\end{proof}

\begin{proof}[Proof of Theorem~\ref{thm:rational}]
Theorem~\ref{thm:normal} supplies the required actual domains for $E_i:=-D_i$ and their adjoints.
Maximal $\Delta$ follows either from \cite[Proposition~3.5]{ACSY24} or Proposition~\ref{prop:bare-maximal}.
Proposition~\ref{prop:intrinsic} proves the equivalence of items~1 and~2.
The fraction algebra just proved gives item~4 from item~1.
Item~4 implies item~3.
Under item~3, closability gives bounded finite-rank extensions; Lemma~\ref{lem:fundamental} and Proposition~\ref{prop:intrinsic} give item~1.
This also proves the realization $T_i=-\overline{\Gamma[Q,D_i]}$ and uniqueness at a fixed fraction.

We check the full natural domains in the statement.
For a bounded representative $K_i$ agreeing on $\M\Omega$, test its adjoint against $w\in\M\Omega$ using \eqref{eq:fraction-adjoint}.
Density identifies that vector with $K_i^*w$.
The same pairing against any $v$ in the natural fraction domain gives $\Gamma[Q,D_i]v=-K_iv$.
Thus the equality holds on the entire natural domain.
Conversely, finite rank there restricts to finite rank on $\M\Omega$.
All these natural domains are dense because they contain $\M\Omega$.
Fraction independence and equality of numerical ranks follow from \eqref{eq:common-defect}.

For a supplied strong system, put $\eta_i:=D_i\Omega$.
The normalized map $-D_i(a\Omega)+a\eta_i$ has adjoint values $-D_i^*(a\Omega)+a\eta_i^*$ on $\M\Omega$, and is therefore closable.
It agrees on the star core with the canonical right slice by the proof of $\mathsf S\Rightarrow\mathsf A$.
The canonical closure is contained in its closure, and both are defined on $\M\Omega$ by Theorem~\ref{thm:normal}; hence there
\[
 D_i^{\rm can}(a\Omega)=D_i(a\Omega)-a\eta_i.
\]
The normalization cancels in each fraction expression because $sb=at$.
Both fraction expressions are closable by the adjoint test, even if their natural domains differ.
If either is finite rank, its bounded closure is determined by the common dense restriction to $\M\Omega$.
The other closed graph contains this everywhere-defined bounded graph and hence equals it.
This proves the final assertion without identifying the two unclosed domains.
\end{proof}

\begin{remark}\label{rem:some-every}
The boundary certificate in Proposition~\ref{prop:intrinsic} quantifies over some fraction.
Maximal $\Delta$ alone does not guarantee a certificate at every prescribed fraction.
For example, let $h$ be uniform on $(0,1)$, let $S$ be a standard semicircular operator free from $h$, and put $Y=hSh$.
The pair $(h,Y)$ has maximal $\Delta$: writing $R=\rho(h)$, $B=\rho(S)$, use
\[
 [RBR,T]=[R,BRT+TRB]+[B,RTR]
\]
to transfer any finite-rank cycle to the free diffuse pair $(h,S)$, which has maximal $\Delta$ by \cite[Section~6]{MSY23}.
The element $S=h^{-1}Yh^{-1}$ is rational.
For its trivial fraction a certificate would force $RT_YR=P_\Omega$ by maximality of $(h,S)$.
Applying this to $\Omega$ would require an $L^2$ vector $v$ with $hv=1$, whereas $h^{-1}\notin L^2(0,1)$.
At the fraction $(h,hS,Sh,h)$, the operators
\[
 T_Y=P_\Omega,\qquad T_h=-P_\Omega Sh-hSP_\Omega
\]
do satisfy the boundary equation.
The same obstruction precludes a full-domain dual at $Y$, since $[E,h]=0$ and $[E,hSh]=P_\Omega$ would give $h(E(Sh)-SE(h))=1$.
Thus this example does not meet the additional hypotheses of Theorem~\ref{thm:rational}.
\end{remark}
